\section{Conclusion}

We study context selection through state-conditioned marginal utility and show
that predictor response provides the key observable for estimating that value.
The response--utility relation admits a controlled expansion for general smooth
losses, an exact Bregman identity, explicit softmax-cross-entropy constants and
a response-sufficiency result under squared loss. These results motivate
response-aware and theory-shaped selectors instead of treating context ranking
as a content-only problem.

The experiments connect that theory to both strong gains and clear regime
changes. When marginal utility is identifiable, response-aware routing
substantially outperforms a broad selector toolbox. Stronger predictors can
retain large oracle opportunity while removing the systematic residual
structure required to rank candidates, and larger candidate libraries can
increase that oracle opportunity while reducing the fraction that routing can
realise. Pool geometry explains when relevance changes sign and yields a
label-free adaptive rule, while a second task family separates the value of
dynamic ranking from the calibration of stopping. The resulting principle is
practical: context selection should be conditioned not only on what information
is available, but also on whether its incremental value is identifiable from
the predictor and deployment state.
